% ECONOMICS_PROBLEM: {"id": "Q54-485", "title": "Measuring climate adaptation", "jel_code": "Q54", "source_id": "OP-0113"}
% !TeX program = xelatex
\documentclass[11pt,a4paper]{article}
\usepackage[margin=23mm]{geometry}
\usepackage{fontspec}
\setmainfont{Latin Modern Roman}
\usepackage{amsmath,amssymb,mathtools}
\usepackage{xcolor,enumitem,booktabs,longtable,array,xurl,hyperref}
\definecolor{muted}{HTML}{53636C}
\hypersetup{colorlinks=true,urlcolor=blue,linkcolor=blue}
\setlength{\parindent}{0pt}
\setlength{\parskip}{5pt}
\newcommand{\R}{\mathbb R}
\newcommand{\E}{\mathbb E}
\newcommand{\N}{\mathbb N}
\newcommand{\ind}{\mathbf 1}
\newcommand{\Var}{\operatorname{Var}}
\newcommand{\Cov}{\operatorname{Cov}}
\newcommand{\supp}{\operatorname{supp}}
\DeclareMathOperator*{\argmin}{arg\,min}
\DeclareMathOperator*{\argmax}{arg\,max}
\newcommand{\entrymeta}[3]{{\sffamily\small\color{muted}\textbf{\texttt{#1}}\quad|\quad #2\par #3\par}}
\newcommand{\Problemref}[1]{\texttt{#1}}
\newcommand{\JELref}[2]{\texttt{#2} (JEL \texttt{#1})}
\begin{document}
\section*{Measuring climate adaptation}
\textbf{Problem ID:} \texttt{Q54-485}\quad \textbf{JEL:} \texttt{Q54}\par
% BEGIN ECONOMICS STATEMENT
\label{problem:OP-0113}
\label{atlas:prob:C3}

\noindent\textbf{Original identifier:} C3 (atlas item 113). \textbf{Field:} Environmental and climate economics. \textbf{Classification:} editorial research family with established restricted solutions; current open status not certified.

\subsection*{Definitions and assumptions}

Let $F$ specify a joint distribution of weather histories over a declared horizon and population. For household or firm $i$, adaptation $a_i\in A_i$ changes outcome $Y_i(w,a_i)$ under weather history $w$ and consumes real resources $k_i(a_i,w)$ measured in a fixed currency year. A welfare index $v_i(Y_i)$ is monetized using a stated marginal consumption value, or is replaced by a clearly defined monetary benefit. Define net welfare
\[
W(F,a)=\sum_i\omega_i\mathbb E_F[v_i(Y_i(w,a_i))-k_i(a_i,w)],
\]
with supplied weights $\omega_i\ge0$. An adaptation regime $a^*(F)$ maximizes $W$ over a stated feasible set, allowing spillovers only if included in $Y_i$ and $k_i$.

\subsection*{Formal research question}

Identify or sharply bound the welfare value of adaptation and the residual climate loss under explicit selection and transport restrictions. For baseline climate $F_0$, changed climate $F_1$ and feasible baseline adaptation $a^0$, define
\[
B_A=W(F_1,a^*(F_1))-W(F_1,a^0),\qquad
L=W(F_0,a^*(F_0))-W(F_1,a^*(F_1)).
\]
Characterize which observational weather panels, adoption interventions and validation data identify these quantities when adaptation and location choices are endogenous. Construct inference over a specified class permitting costly defensive spending, heterogeneous exposure and bounded unobserved selection, and separate the sampling error from ambiguity in long-run climate responses and adaptation costs.

\subsection*{Interpretation and scope}

Reduced temperature sensitivity of mortality or yields can coexist with large resource costs, changed composition or displaced damages. Thus a change in an outcome coefficient is neither the welfare benefit $B_A$ nor proof of costless adaptation. If $a^0$ is feasible and an optimum exists, $B_A\ge0$ by definition; observed adoption need not be optimal. Weather fluctuations identify responses within sampled weather and institutional regimes, while a shift in $F$ can change technologies, prices and expectations. Migration changes the population at risk; holding original residents fixed and evaluating future residents are different targets. Transfers, medical expenditures and physical resource losses must be handled consistently to avoid double counting.

\subsection*{Source audit and status}

[S1] examines weather, mortality and residential energy use; [S2] studies historical declines in temperature mortality, including adaptation; [S3] estimates agricultural adaptation using US climate trends. These verified empirical sources already estimate selected adaptation channels. They do not identify every net-welfare component or universal future response. The selection-aware full-welfare target is an editorial specialization, and the bibliography does not certify its precise 2026 open status.

\subsection*{References}

\begin{enumerate}[label={[\textnormal{S}\arabic*]},leftmargin=*]

\item Olivier Deschênes and Michael Greenstone (2011). \emph{Climate Change, Mortality, and Adaptation: Evidence from Annual Fluctuations in Weather in the US}. American Economic Journal: Applied Economics 3(4), 152–185. \url{https://doi.org/10.1257/app.3.4.152}. \textbf{Locator:} Publisher or institutional abstract and bibliographic record. \textbf{Support and limits:} Weather mortality and residential energy responses; weather variation does not alone identify long-run adaptive welfare.

\item Alan Barreca, Karen Clay, Olivier Deschenes, Michael Greenstone and Joseph S. Shapiro (2016). \emph{Adapting to Climate Change: The Remarkable Decline in the US Temperature-Mortality Relationship over the Twentieth Century}. Journal of Political Economy 124(1), 105–159. \url{https://doi.org/10.1086/684582}. \textbf{Locator:} Publisher or institutional abstract and bibliographic record. \textbf{Support and limits:} Historical temperature-mortality adaptation, including air conditioning; not all adaptation channels or future locations.

\item Marshall Burke and Kyle Emerick (2016). \emph{Adaptation to Climate Change: Evidence from US Agriculture}. American Economic Journal: Economic Policy 8(3), 106–140. \url{https://doi.org/10.1257/pol.20130025}. \textbf{Locator:} Publisher or institutional abstract and bibliographic record. \textbf{Support and limits:} US agriculture estimates comparing longer-run climate trends; transport requires extra assumptions.

\end{enumerate}

\subsection*{Common conventions: Source mathematical conventions}

Unless an entry states otherwise, agent and firm sets are finite. State and action spaces are measurable, strategies and outcome maps are measurable, probability laws are specified, and expectations exist for the targets under discussion. Infinite-horizon discount factors lie in $(0,1)$ and discounted objectives are finite. Norms, tolerances and complexity input sizes must be specified for computational comparisons. Constants or primitives that appear locally are inputs, not empirical facts asserted by this catalogue.

\textbf{Identification.} A statistical model $m\in\mathcal M$ implies an observable law $P_m$ and target $\theta(m)$. At law $P$, its identified set is
\[
\Theta(P;\mathcal M)=\{\theta(m):m\in\mathcal M,\ P_m=P\}.
\]
Point identification holds when this set is a singleton, or equivalently when $P_{m_1}=P_{m_2}$ implies $\theta(m_1)=\theta(m_2)$ for all compatible models. Sharp bounds mean bounds attained, or approached when the boundary is not attained, by compatible models. Writing this set is a definition, not a solution: the substantive task is to establish informative restrictions and derive or estimate the set. An unrestricted-confounding model may give only trivial bounds.

\textbf{Inference.} Identification, estimability and valid uncertainty quantification are separate questions. Fix $\alpha\in(0,1)$. A confidence procedure $C_n$ over a declared law class $\mathcal P$ has asymptotic uniform coverage at level $1-\alpha$ when
\[
\liminf_{n\to\infty}\inf_{P\in\mathcal P}\Pr_P\{\theta(P)\in C_n\}\geq1-\alpha.
\]
For set-identified targets the coverage event must be defined separately, for example coverage of the full identified set. If an entry uses identification-indexed models instead of uniquely indexed laws, the same coverage convention is applied over those models. Pointwise limits do not establish uniform validity, and asymptotic validity does not guarantee finite-sample precision. Sample sizes, dependence structures and weak-identification sequences are part of each application.

\textbf{Counterfactuals.} $Y_i(a)$ is a potential outcome under a specified intervention, including its timing, implementation and versions. It is not $Y_i$ conditional on voluntarily choosing $a$. A full intervention vector permits interference; shorthand scalar treatment notation does not silently impose no interference. Assignment, selection, attrition, anticipation and measurement must be specified. A claim about an instrument's local effect is not automatically a population policy effect.

\textbf{Economic equilibrium.} A counterfactual model includes best responses, constraints, feasibility, market clearing, state transitions and expectations consistent with the stated information. When several equilibria exist, either a declared selector is an input or the target is set-valued. Comparisons across policy regimes must use the stated selection convention. Reduced-form relationships are not presumed invariant after an equilibrium-changing intervention.

\textbf{Optimization and welfare.} Optimization is over the stated feasible set. If attainment is not established, $\sup$ or $\inf$ and $\varepsilon$-optimal choices replace an asserted maximizer or minimizer. For example, with value $V=\sup_{a\in\mathcal A}W(a)<\infty$, an $\varepsilon$-optimal set is $\{a\in\mathcal A:W(a)\geq V-\varepsilon\}$ for $\varepsilon>0$. Benefits, costs, risk and health or environmental outcomes can be aggregated only using declared units and conversion weights. Interpersonal utility weights, the treatment of fiscal transfers and foreign profits, discounting, and the behavioral welfare criterion are explicit inputs. A comparative-static derivative additionally requires existence and appropriate differentiability of the selected counterfactual.

\textbf{Sources.} Local labels [S1], [S2], and so forth restart in every question. Locators identify the relevant abstract, model, section or result at the level actually checked; a bibliographic or abstract-level verification is not represented as inspection of an entire proof. Working papers are distinguished from later publications and changing versions. Source summaries are paraphrased. All URL checks and editorial formulations refer to the audit date stated above.
% END ECONOMICS STATEMENT
\end{document}
